\chapter{Fase 1 -- Interpolação horizontal nativa}
\label{cap:fase1}

\section{Objetivo e reaproveitamento de código}
\label{sec:fase1_reuso}

O objetivo da Fase~1 é remapear, célula a célula, os campos escalares já
extraídos e verticalmente interpolados para níveis de pressão fixa
(estágio herdado sem modificação da rota original, programa
\texttt{extract\_fields}) da malha nativa \emph{global} de origem para os
centros de célula da malha nativa \emph{regional} de destino, aplicando
a interpolação baricêntrica descrita na \cref{sec:baricentrica}.

A decisão de engenharia central desta fase foi \emph{não reescrever} o
motor de cálculo de pesos baricêntricos e localização de ponto no
triângulo dual: o utilitário \texttt{convert\_mpas} (mantido como
dependência da rota original, propriedade da NCAR) já implementa esse
mecanismo internamente -- confirmado por leitura de código-fonte,
seção ``Interpolation methodology'' de sua documentação -- para o caso
de remapear \emph{para} uma grade latitude-longitude regular. A
investigação encontrou que o módulo interno responsável por montar a
malha de destino (\texttt{target\_mesh.F90}) já possui um modo
alternativo, não exposto pela interface de linha de comando original,
que aceita uma \emph{lista de pontos dispersos} (\emph{scatter}) em vez
de uma grade regular -- bastando fornecer os vetores de latitude e
longitude dos pontos de destino com \texttt{irank=1}. O motor de cálculo
de pesos propriamente dito (\texttt{remapper.F90}) é inteiramente
agnóstico entre os dois modos.

O programa \texttt{hinterp\_native.F90}, novo neste trabalho, é
essencialmente uma cópia da estrutura de \texttt{convert\_mpas.F90} (os
mesmos módulos de leitura de malha, varredura de campos e escrita de
saída), com uma única mudança: em vez de chamar a rotina de montagem de
grade regular, ele lê as coordenadas \texttt{latCell}/\texttt{lonCell}
da malha de destino e as passa ao módulo de montagem de malha em modo
\emph{scatter}. Nenhuma linha do motor de pesos baricêntricos
(\texttt{remapper.F90}) foi modificada.

\section{Procedimento}

Para cada campo escalar $\phi$ requisitado (temperatura, componentes de
vento, altura geopotencial, umidade específica e relativa, e as
respectivas variáveis de superfície), e para cada célula $\mathbf{p}_j$
da malha de destino:

\begin{enumerate}
  \item Localiza-se o triângulo de Delaunay da malha de origem que
    contém $\mathbf{p}_j$, usando a conectividade célula-célula/vértice
    já armazenada na malha de origem (o dual de Delaunay não precisa ser
    recalculado do zero -- é uma propriedade geométrica derivável
    diretamente de \texttt{cellsOnVertex}/\texttt{verticesOnCell}, já
    presentes no arquivo de malha).
  \item Calculam-se os três pesos baricêntricos $(\lambda_A, \lambda_B,
    \lambda_C)$ via \cref{eq:baricentrica_area}, com as áreas computadas
    sobre a superfície esférica.
  \item Aplica-se \cref{eq:baricentrica} para obter
    $\phi(\mathbf{p}_j)$.
\end{enumerate}

O resultado é escrito em um arquivo NetCDF intermediário
(\texttt{native\_target.nc}), com os campos ainda organizados por nível
de pressão fixa mas já localizados nos centros de célula da malha
regional -- pronto para a interpolação vertical da Fase~3
(\cref{cap:fase3}).

\begin{caixaachado}
Um efeito colateral do reúso do módulo \texttt{remapper.F90} sem
modificação foi que suas rotinas internas de geração das variáveis de
coordenada de saída (\texttt{remap\_get\_target\_latitudes}/
\texttt{longitudes}) assumem semântica de grade regular e, em modo
\emph{scatter}, colapsam a variável \texttt{latitude} de saída para um
único valor escalar em vez de um vetor por ponto -- um bug de uso do
módulo original (não corrigido no \texttt{convert\_mpas} em si) que
passou despercebido por não ser exercitado no caminho de código
tradicional (sempre grade). A correção adotada foi \emph{aditiva}, não
substitutiva: em vez de reescrever essas rotinas (uma primeira tentativa
nesse sentido quebrou a escrita do NetCDF, com erro de ordem de criação
de dimensões em modo \texttt{define}), o programa novo escreve campos
adicionais corretos, \texttt{target\_latitude}/\texttt{target\_longitude},
por cima da saída original -- preservando as chamadas herdadas intactas.
\end{caixaachado}

\section{Reconstrução do vento nas arestas}
\label{sec:vento_arestas}

O vento não é tratado como um campo escalar simples nesta fase, pelo
motivo estrutural já discutido na \cref{sec:dualidade}: sua componente
prognóstica no \mpasa{} é a componente \emph{normal} $u$ armazenada nas
arestas entre células, não um vetor $(u,v)$ no centro de célula. Seguindo
a evidência apresentada por Ha et al.~\cite{ha2017enkf} de que reconstruir
o vento a partir das componentes zonal/meridional originais do
first-guess (que já estão disponíveis, por definição, nos centros de
célula da malha de \emph{origem}) e só então projetar na direção normal
da aresta de \emph{destino} produz resultado superior a interpolar
diretamente uma quantidade já projetada, o procedimento adotado -- ainda
na etapa combinada com a Fase~3 (\cref{cap:fase3}) -- calcula, para cada
aresta $e$ da malha de destino com células vizinhas $c_1(e), c_2(e)$:
\begin{equation}
  u_{fg}(e) = \cos(\alpha_e)\, \bar{u}_{zonal} + \sin(\alpha_e)\, \bar{v}_{meridional},
  \qquad
  \bar{(\cdot)} = \tfrac{1}{2}\big[(\cdot)_{c_1(e)} + (\cdot)_{c_2(e)}\big],
  \label{eq:vento_projecao}
\end{equation}
onde $\alpha_e$ é o ângulo da aresta $e$ da malha de \emph{destino} em
relação ao meridiano local (variável \texttt{angleEdge}, já presente na
malha) e as médias $\bar{u}_{zonal}$, $\bar{v}_{meridional}$ são tomadas
sobre os valores \emph{já remapeados horizontalmente} (via
\cref{eq:baricentrica}, campo a campo) nos dois centros de célula
vizinhos da aresta de destino. Esse perfil vertical projetado é então
interpolado verticalmente (\cref{cap:fase3}) usando como coordenada de
altura a média das alturas geopotenciais das duas células vizinhas.

\section{Validação numérica}
\label{sec:validacao_fase1}

A validação desta fase explora diretamente a propriedade de exatidão nos
vértices discutida na \cref{sec:baricentrica}: interpolando os campos de
uma malha de teste (\texttt{x1.2562.static.nc}, 2562 células) para os
próprios centros de célula dessa mesma malha (\emph{auto-\emph{scatter}}),
o erro obtido foi \emph{identicamente nulo} em todas as 2562 células,
tanto nos valores de campo quanto nas coordenadas de saída
(\texttt{target\_latitude}/\texttt{target\_longitude}). Como comparação
de controle, o mesmo experimento repetido passando por uma grade lat-lon
intermediária fina ($\SI{0.25}{\degree}$) -- reproduzindo o caminho da
rota original -- introduziu um erro sistemático mensurável, da ordem de
$10^{-4}$ a $10^{-3}$ (em unidades do campo), \emph{mesmo nesses mesmos
pontos de coincidência exata}, confirmando de forma independente que o
erro de ida-e-volta discutido na \cref{cap:introducao} é real e
mensurável, não apenas uma preocupação teórica.
